\section{Related Work}
\label{sec:related}

\id{-nya} is analysed as a polysemous Indonesian clitic with possessive, definite, anaphoric, and nominalising uses~\cite{denistia2021morphology}. Computational resources such as MorphInd~\cite{morphind_larasati}, IndoMorph~\cite{indomorph2025}, and Larasati's clitic work~\cite{larasati2012clitics} show that Indonesian cliticisation is a recognised modelling problem, but they do not evaluate retrieval effectiveness.

Indonesian IR work has mostly studied stemming rather than clitic-specific preprocessing. Tala~\cite{tala2003stemming} found stemming useful in a pre-neural IR setting, and later work compares stemming accuracy for systems such as Sastrawi~\cite{sastrawi} and Nazief--Adriani~\cite{nazief1996indonesian}. Modern Indonesian and multilingual resources, including IndoNLU~\cite{wilie2020indonlu}, IndoLEM~\cite{koto2020indolem}, NusaCrowd~\cite{cahyawijaya2023nusacrowd}, MIRACL~\cite{zhang2023miracl}, and Mr.\ TyDi~\cite{zhang2021mrtydi}, enable stronger evaluation, but none isolates \id{-nya} preprocessing.

Prior work suggests that anaphora resolution can improve retrieval or retrieval-augmented generation in English~\cite{mitkov2002anaphora,jang2025coref}. Indonesian coreference systems exist, including a CNN mention-pair resolver~\cite{auliarachman2020coref} and IndoCoref~\cite{artari2021indocoref}, but they are not plug-in raw-passage retrieval preprocessors. We therefore test an in-house deterministic surface heuristic, while recognising that stronger coreference models remain future work. More broadly, surveys of Transformer preprocessing report mixed, task-dependent effects from classical preprocessing~\cite{camacho2023preprocessing}, motivating a retrieval-specific Indonesian test.
